% Generated draft: requires mathematical and source review.
\documentclass[11pt]{article}
\usepackage[margin=25mm]{geometry}
\usepackage{fontspec}
\setmainfont{Latin Modern Roman}
\usepackage{amsmath,amssymb,mathtools,unicode-math}
\setmathfont{Latin Modern Math}
\usepackage{enumitem,needspace,xurl,hyperref}
\setlength{\parindent}{0pt}
\setlength{\parskip}{5pt}
\setlength{\emergencystretch}{3em}
\newcommand{\N}{\mathbb{N}}
\newcommand{\Z}{\mathbb{Z}}
\newcommand{\Q}{\mathbb{Q}}
\newcommand{\R}{\mathbb{R}}
\newcommand{\C}{\mathbb{C}}
\newcommand{\F}{\mathbb{F}}
\newcommand{\A}{\mathbb{A}}
\newcommand{\PP}{\mathbb{P}}
\newcommand{\E}{\mathbb{E}}
\newcommand{\eps}{\varepsilon}
\newcommand{\dd}{\,\mathrm{d}}
\begin{document}
\section*{UnsolvedMath ID: 700723}
\textbf{Source problem code:} catalogue entry
\subsection{Definitions and mathematical statement}
A group is torsion-free if its identity is its only element of finite order. A group \(G\) is soluble if its derived series, defined by \(G^{(0)}=G\) and \(G^{(i+1)}=[G^{(i)},G^{(i)}]\), reaches the trivial group after finitely many steps.

For the following working formulation, interpret rank as Prüfer rank. For a finitely generated group \(K\), let \(d(K)\) denote the minimum size of a generating set, with \(d(\{1\})=0\). Define
\[
\operatorname{rk}(G)=\sup\{d(K):K\leq G\text{ is finitely generated}\}\in\Z_{\geq0}\cup\{\infty\},
\]
where an unbounded supremum is denoted by \(\infty\). The supplied record does not define rank, so this interpretation requires confirmation against the original source.

The question attributed to D. I. Zaitsev is whether
\[
\forall G\ \bigl(G\text{ torsion-free and soluble}\bigr)\quad
\forall H\leq G\ \bigl([G:H]<\infty\bigr),\qquad
\operatorname{rk}(G)=\operatorname{rk}(H).
\]
No normality assumption on \(H\) is imposed, and no finite-rank hypothesis is imposed on \(G\).

For the contextual qualifications, use the conventional meaning of a rational series: a finite subnormal series
\[
\{1\}=G_0\triangleleft G_1\triangleleft\cdots\triangleleft G_m=G
\]
whose factors \(G_i/G_{i-1}\) are isomorphic to nonzero additive subgroups of \(\Q\). The trivial group may have the empty series. The supplied record does not specify this convention, so it too requires source confirmation.

The source reports an affirmative answer for torsion-free soluble groups having a rational series, citing Zaitsev's 1971 contribution. It also reports that every torsion-free soluble group of finite rank contains a subgroup of finite index having a rational series. These are retained as supplied contextual assertions, not independently verified results. The local registry labels the problem open; that label is supplied metadata rather than an independently verified assessment.
\subsection{Short English statement}
Does every finite-index subgroup of a torsion-free soluble group have the same rank as the whole group? Here rank is provisionally interpreted as Prüfer rank. The source reports a positive answer for groups with a rational series and states that every finite-rank torsion-free soluble group has a finite-index subgroup with such a series.
\subsection{Sources}
\begin{itemize}
\item Local UnsolvedMath registry, permanent ID \(700723\), exact source problem code \texttt{catalogue entry}, title \emph{Kourovka Notebook Problem 11.30}. The supplied record attributes the problem to D. I. Zaitsev, gives the proposed year as 1990, and labels its status \texttt{open}.
\item \emph{The Kourovka Notebook (Unsolved Problems in Group Theory)}, Problem 11.30, originating in the 11th issue (1990). The supplied background identifies the consulted edition as the 21st issue edition (2026), \texttt{arXiv:1401.0300v48}. The registry collection describes Issues 1--21 as edited by E. I. Khukhro and V. D. Mazurov. These bibliographic details are supplied attribution and have not been independently checked.
\item D. I. Zaitsev, contribution in \emph{Groups with restrictions on subgroups}, Naukova dumka, Kiev, 1971, pp. 115--130 (Russian). The supplied statement cites this contribution for the affirmative answer for groups having a rational series; it supplies no contribution title.
\end{itemize}
\end{document}
